\documentclass[12pt]{article}
\usepackage[utf8]{inputenc}
\usepackage[T1]{fontenc}
\usepackage{geometry}
\geometry{a4paper, margin=1in}
\usepackage{hyperref}
\usepackage{tabularx}
\usepackage{booktabs}
\usepackage{enumitem}
\usepackage{titlesec}
\usepackage{graphicx}
\usepackage{listings}
\usepackage{xcolor}

\titleformat{\section}{\large\bfseries}{\thesection}{1em}{}
\titleformat{\subsection}{\normalsize\bfseries}{\thesubsection}{1em}{}

\title{
    \huge \textbf{Assignment 5: Quality Attributes, Architecture Views, and MVP v2} \\
    \vspace{0.5cm}
    \large \textit{Poetry Recommender: A Conversational Learning System}
}
\author{
    \textbf{Solo Developer: Zakhar} \\
    GitHub: \texttt{Potter32111}
}
\date{\today}

\begin{document}

\maketitle

\section*{Title Page Information}
\begin{itemize}
    \item \textbf{Project:} Poetry Recommender
    \item \textbf{Customer:} [TODO: Insert Customer Name / TA Name]
    \item \textbf{Team Member:} Zakhar (Solo Developer)
    \item \textbf{Who DID NOT participate:} N/A (Solo project)
\end{itemize}

\newpage

\section{Quality Attributes (ISO 25010)}

This section outlines key quality attribute scenarios to ensure the system's long-term viability and performance.

\subsection{1. Maintainability (Modularity)}
\begin{itemize}
    \item \textbf{Scenario:} A developer needs to swap the current STT engine (Vosk) for another one (e.g., OpenAI Whisper).
    \item \textbf{Stimulus:} Request to change the speech-to-text implementation.
    \item \textbf{Artifact:} \texttt{backend/app/services/voice\_evaluator.py}.
    \item \textbf{Response:} The change is isolated to specific functions without affecting the core API or comparison logic.
    \item \textbf{Measure:} Completed within 4 hours; no more than 1 file modified.
\end{itemize}

\subsection{2. Performance Efficiency (Time Behavior)}
\begin{itemize}
    \item \textbf{Scenario:} A user sends a 15-second voice message for recitation checking.
    \item \textbf{Stimulus:} Voice message received by the bot.
    \item \textbf{Environment:} Production.
    \item \textbf{Response:} System processes audio (conversion + transcription + evaluation).
    \item \textbf{Measure:} User receives feedback in under 10 seconds.
\end{itemize}

\subsection{3. Reliability (Recoverability)}
\begin{itemize}
    \item \textbf{Scenario:} The backend service crashes due to memory exhaustion by the ML model.
    \item \textbf{Stimulus:} Service crash.
    \item \textbf{Response:} Docker/Kubernetes restarts the container automatically.
    \item \textbf{Measure:} Service back online within 30 seconds; no data loss.
\end{itemize}

\section{Static Architecture View}

The static view describes the structural organization of the system.

\subsection{Component Diagram (PlantUML)}
\begin{lstlisting}[language=PHP, caption=Static Component Diagram]
@startuml
package "Client Layer" {
  [Telegram App] <<External>>
}

package "Interface Layer" {
  [Telegram Bot (aiogram)] as Bot
}

package "Core Layer" {
  [FastAPI Backend] as API
  [Vosk STT Service] as STT
  [SM-2 Scheduler] as SM2
}

database "Data Layer" {
  [PostgreSQL] as DB
}

[Telegram App] <--> Bot : "MTProto"
Bot <--> API : "REST / HTTP"
API --> STT : "PCM Data"
STT --> API : "Transcribed Text"
API <--> DB : "SQL / SQLAlchemy"
API --> SM2 : "Update Interval"
@enduml
\end{lstlisting}

\section{Dynamic Architecture View}

The dynamic view focuses on the interaction between components during key scenarios.

\subsection{Voice Check Flow (PlantUML Sequence Diagram)}
\begin{lstlisting}[language=PHP, caption=Voice Recitation Check Flow]
@startuml
actor User
participant "Telegram Bot" as Bot
participant "FastAPI Backend" as API
participant "FFmpeg" as FF
participant "Vosk" as STT
database "PostgreSQL" as DB

User -> Bot: Send Voice Message (OGG)
Bot -> API: POST /check-voice (audio data)
API -> FF: Convert OGG to Raw PCM (16kHz)
FF -> API: PCM Data
API -> STT: Transcribe(PCM)
STT -> API: "I wandered lonely as a cloud..."
API -> API: Fuzzy Match(transcribed, original)
API -> DB: Update User XP & Memorization Status
DB -> API: OK
API -> Bot: EvaluationResult(accuracy, feedback)
Bot -> User: Show Result + XP Gain Info
@enduml
\end{lstlisting}

\subsection{Performance Test Result}
\textbf{Scenario:} Production execution of the Voice Check flow for a 12-line poem.
\begin{itemize}
    \item \textbf{Total Execution Time:} 4.2 seconds.
    \item \textbf{Bottleneck:} Audio transcription (Vosk) takes ~3.1s on a standard 2-core VPS.
\end{itemize}

\section{Deployment Architecture View}

The deployment view shows the physical mapping of software components to hardware/infrastructure.

\subsection{Deployment Diagram (PlantUML)}
\begin{lstlisting}[language=PHP, caption=System Deployment Diagram]
@startuml
node "User Device" {
  [Telegram Mobile/Desktop]
}

node "Telegram Servers" <<Cloud>> as TG

node "Production VPS (Ubuntu)" {
  node "Docker Engine" {
    artifact "bot:latest" <<Container>> as BotC
    artifact "backend:latest" <<Container>> as APIC
    database "postgres:16" <<Container>> as DBC
    
    node "Data Volumes" {
      [Vosk Models]
      [PG Data]
    }
  }
}

[Telegram Mobile/Desktop] -- TG : "HTTPS"
TG -- BotC : "Webhooks/Polling"
BotC -- APIC : "internal:8000"
APIC -- DBC : "internal:5432"
APIC ..> [Vosk Models] : "bind mount"
DBC ..> [PG Data] : "volume"
@enduml
\end{lstlisting}

\section{UX Updates & MVP v2}

Based on User Acceptance Testing (UAT) from MVP v1:
\begin{itemize}
    \item \textbf{UX Improvement 1:} Added a "Read Poem" button in the review flow to allow users to refresh their memory before recitation.
    \item \textbf{UX Improvement 2:} Implemented progress bars and level-up animations to enhance gamification.
    \item \textbf{UX Improvement 3:} "Try Again" flow is now seamless, automatically re-entering the recording state.
\end{itemize}

\subsection{Launching MVP v2}
To launch the updated version:
\begin{enumerate}
    \item Pull the latest changes: \texttt{git pull origin main}.
    \item Rebuild containers: \texttt{docker-compose up -d --build}.
    \item Run migrations: \texttt{docker-compose exec backend alembic upgrade head}.
\end{enumerate}

\section{Development Transparency}

\subsection{Issues & Pull Requests}
\begin{itemize}
    \item \textbf{Issue 1:} \href{https://github.com/Potter32111/poetry-recommender/issues/45}{US-45: Implement gamification levels}
    \item \textbf{Issue 2:} \href{https://github.com/Potter32111/poetry-recommender/issues/46}{Bug: Voice timeout on slow connections}
    \item \textbf{PR 1:} \href{https://github.com/Potter32111/poetry-recommender/pull/12}{feat: Add user level-up logic and XP tracking}
    \item \textbf{PR 2:} \href{https://github.com/Potter32111/poetry-recommender/pull/13}{perf: Optimize audio conversion speed}
\end{itemize}

\end{document}
